\scanpage{34}
\mspage{21}

La démonstration est \inserted{essentiellement} la $=$ que pour le cas particulier déjà traité $I'=\varnothing,\ I''=I$, et on aurait pu aussi bien rédiger la démonstration dans ce cadre, à l’aide d’une suite \inserted{str. croissante,} transfinie d’ouverts \inserted{de $I$} $I_\alpha$, avec $\boxed{I_0=I'}$, $I_{\alpha_0}=I''(=I)$.

\deleted{On peut d’ailleurs ramener le cas général au cas déjà traité, en introduisant, quand $I'\subset I(=I'')$ \unsure{fixés|fixé}, \deleted{\unsure{des}} une nouvelle flèche
\[
 u_\bullet:X'_\bullet\longrightarrow Y'_\bullet,
\]
en posant
\[
 X'_i=\underset{j\in I'\cup I_{/i}}{\underrightarrow{\lim}}\,X_j,\quad
 Y'_i=\underset{j\in I'\cup I_{/i}}{\underrightarrow{\lim}}\,Y_j
\]
et \deleted{en travaillant dans} $\unsure{\mathcal M|\mathscr M}'={}_{\unsure{x|X}_{I'}}$
\[
 X'_{\varnothing}=X_{I'},\quad Y'_{\varnothing}=Y^{\unsure{4|\leftarrow}}_{I'}
\]
$\unsure{X'}$}

et

\[
\begin{aligned}
 I_\alpha&=I_{\alpha-1}\cup I_{/s_{\alpha-1}}
 &&\text{si }\alpha\text{ non limite, }\inserted{I_{\alpha-1}\ne I,}\ \deleted{\text{pour}}\ \text{choisissant}\\
 &&&s_{\alpha-1}\in I\smallsetminus I_{\alpha-1}\\
 I_\alpha&=\bigcup_{\beta<\alpha} I_\beta\quad .
 &&\text{si }\alpha\text{ limite.}
\end{aligned}
\]

et posant encore

\[
 X_\alpha=X_{I_\alpha},\quad Y_\alpha=Y_{I_\alpha},\quad u_\alpha=u_{I_\alpha}:X_\alpha\mathbin{\unsure{-|\to}}Y_\alpha.
\]

Soit $\alpha_0$ tel que \deleted{$\unsure{\alpha}$} $I_{\alpha_0}=I$ (indice terminal de la suite) ; \deleted{\unsure{alors} \unsure{il} \unsure{prouv.}} il faut prouver que \unsure{dans|les|la}

\diagram{\begin{tikzcd}[ampersand replacement=\&]
 X_0\arrow[r]\arrow[d,no head]\&X_{\alpha_0}\arrow[d,no head]\\
 Y_0\arrow[r]\&Y_{0,\alpha_0}\arrow[d,"{u_{0,\alpha_0}}"]\\
 \&Y_{\alpha_0}
\end{tikzcd}}
% archive sheet number: 33

